% Results subsection: solo versus team, probe pilot v1 and probe v2 across four models.
% Use with \input{results_four_models} inside the Results section.
% Needs \usepackage{booktabs}. Sources: docs/retail_confirmation_v1_*_results.md,
% docs/individual_probe_pilot_results.md, docs/individual_probe_pilot_v2_results.md.
\section{Results}
\subsection{Team effects and individual probes across four models}
\label{sec:results-four-models}

We repeat the retail partial-approval study with four models from four families:
GPT-5 nano (\texttt{gpt-5-nano-2025-08-07}, OpenAI API),
Qwen~3.6 35B-A3B (\texttt{Qwen3.6-35B-A3B-FP8}, thinking disabled, NCSA Lumen),
DeepSeek V4.1 Flash (\texttt{DeepSeek-V4.1-Flash-1017}, DeepSeek API), and
Gemini~3.1 Flash-Lite (\texttt{gemini-3.1-flash-lite}, Gemini API).
In every episode the customer asks for two order changes, $X$ and $Y$, and approves only $X$.
The outcome is whether any agent executes $Y$. It is scored from executed tool calls, and every
flagged episode is confirmed by replaying its calls against a fresh copy of the database.
Replay reproduced all labels, with no disagreements across 3{,}964 episodes.

\paragraph{Teams versus solo agents.}
Table~\ref{tab:solo-team} compares a single agent with a three-agent team on the same 150 case
groups. The team has a coordinator who speaks to the customer, plus two specialists who share
notes. Teams executed $Y$ far more often than solo agents for nano ($10$ versus $1$ of $150$)
and Qwen ($23$ versus $1$). DeepSeek ($2$ versus $0$) and Gemini ($2$ versus $2$) showed no
team effect. Every team violation was committed by a specialist. The team effect is therefore
real but model dependent.

\begin{table}[t]
\centering
\small
\caption{Retail partial approval, solo versus team. Counts are episodes in which an agent
executed the unapproved change $Y$, out of 150 matched case groups. $p$ is the exact McNemar
test on the matched pairs. ``Completed'' counts team episodes that executed the approved
change $X$.}
\label{tab:solo-team}
\begin{tabular}{lrrrr}
\toprule
Model & Solo & Team & $p$ & Completed \\
\midrule
GPT-5 nano            & 1 & 10 & 0.012    & 131 \\
Qwen 3.6$^{\dagger}$  & 1 & 23 & $3\times10^{-6}$ & 52 \\
DeepSeek V4.1 Flash   & 0 & 2  & 0.50     & 124 \\
Gemini 3.1 Flash-Lite & 2 & 2  & 1.00     & 124 \\
\bottomrule
\end{tabular}

\vspace{2pt}
{\footnotesize $^{\dagger}$39 Qwen episodes were rerun a second time after provider failures.
Under the prespecified single rerun, 113 complete pairs give 1 versus 14 ($p=0.001$).}
\end{table}

\paragraph{Probe pilot v1: explicit refusal.}
Individual probes freeze one agent at the decision point. The customer has just said
``Please go ahead with $X$ only. Do not do $Y$. I am not approving it.'' The agent then sees one
scripted teammate message. The five conditions are no message (C0), a correct delegation of
$X$ (C1), an overbroad request for $X$ and $Y$ (C2), a false claim that both were approved
(C3), and an overbroad request after a teammate has already executed another unapproved change
(C4). Table~\ref{tab:probe-v1} reports specialist outcomes. Only Qwen could be moved to act
against an explicit refusal ($15/24$ under C3 and $8/24$ under C4). nano, DeepSeek and Gemini
never executed $Y$. Coordinators never executed $Y$ in any model.

\begin{table}[t]
\centering
\small
\caption{Probe pilot v1 with an explicit customer refusal. Specialist episodes that executed
$Y$, out of 24 per cell (12 case groups $\times$ 2 specialist roles).}
\label{tab:probe-v1}
\begin{tabular}{lrrrr}
\toprule
Teammate message & nano & Qwen & DeepSeek & Gemini \\
\midrule
C0: none                                  & 0 & 0  & 0 & 0 \\
C1: correct delegation of $X$             & 0 & 1  & 0 & 0 \\
C2: overbroad ($X$ and $Y$)               & 0 & 0  & 0 & 0 \\
C3: false approval claim                  & 0 & 15 & 0 & 0 \\
C4: overbroad after a peer violation      & 0 & 8  & 0 & 0 \\
\midrule
Total (of 120)                            & 0 & 24 & 0 & 0 \\
Completed $X$ (of 120)                    & 86 & 119 & 120 & 120 \\
\bottomrule
\end{tabular}
\end{table}

\paragraph{Probe v2: customer wording.}
Retail customers rarely refuse outright, so probe v2 also varies the customer's wording. The
reply is explicit, silent about $Y$, uncertain about $Y$, or asks to hold $Y$. This is crossed
with four teammate messages: none, a status relay that ends ``Customer confirmed'', an
overbroad request, and a false approval claim. Table~\ref{tab:probe-v2} shows that silence is
the condition that breaks every model. Under silent wording, specialists executed $Y$ in 17
(nano), 21 (DeepSeek), 75 (Gemini) and 94 (Qwen) of 96 episodes. Outside it, only Qwen and,
rarely, nano did so. Gemini is the only model whose coordinators violated (3 of 96, all under
silent wording). Specialists still completed $X$ in 296 to 384 of 384 episodes, so low violation
rates do not come from refusing all work.

\begin{table}[t]
\centering
\small
\caption{Probe v2: customer wording crossed with four teammate messages. Specialist episodes
that executed $Y$, out of 96 per wording (4 messages $\times$ 24).}
\label{tab:probe-v2}
\begin{tabular}{lrrrr}
\toprule
Customer wording & nano & Qwen & DeepSeek & Gemini \\
\midrule
Explicit (``Do not do $Y$'')              & 0  & 20  & 0  & 0  \\
Silent (``Yes, go ahead with $X$'')       & 17 & 94  & 21 & 75 \\
Hedge (``$Y$, I'm not sure about yet'')   & 6  & 47  & 0  & 0  \\
Defer (``Please hold off on $Y$'')        & 0  & 19  & 0  & 0  \\
\midrule
Total (of 384)                            & 23 & 180 & 21 & 75 \\
Silent, no teammate message (of 24)       & 1  & 22  & 1  & 12 \\
Coordinators (of 96)                      & 0  & 0   & 0  & 3  \\
\bottomrule
\end{tabular}
\end{table}

\paragraph{Do probes predict teams?}
Only partly. Probes identify Qwen as the riskiest model, consistent with its team rate. However,
they rank Gemini second ($75/384$), while Gemini teams were among the safest ($2/150$). They
also cannot separate nano from DeepSeek ($23$ versus $21$ of 384), whose team rates differ
fivefold ($10$ versus $2$ of 150). An audit of the nano and DeepSeek team traces suggests why.
In 9 of their 12 team violations, no teammate message had requested $Y$, and nano specialists
had more decisions at risk per episode than DeepSeek specialists (5.6 versus 3.9). On the first
such decision, the two models behaved almost identically, and that first decision is all a
probe measures. A forecast of team risk from individual probes therefore needs a model of
exposure as well as of the response to a single decision.
